\documentclass[10pt,journal,compsoc]{IEEEtran}

\usepackage{amsmath,amssymb,amsfonts}
\usepackage{algorithmic}
\usepackage{graphicx}
\usepackage{textcomp}
\usepackage{xcolor}
\usepackage{booktabs}
\usepackage{microtype}
\usepackage{hyperref}
\usepackage{cite}

\hypersetup{
    colorlinks=true,
    linkcolor=blue,
    filecolor=magenta,      
    urlcolor=cyan,
    citecolor=blue,
}

\begin{document}

\title{Real-Time On-Device Vehicle Sheet-Metal Segmentation and Dichromatic Paint Relighting for Mobile Augmented Reality}

\author{Anonymous~Authors%
\thanks{Manuscript received October 2026; revised January 2027. Corresponding author: Academic Research Group.}}

\markboth{IEEE Transactions on Visualization and Computer Graphics,~Vol.~XX, No.~Y, October~2026}%
{Anonymous \MakeLowercase{\textit{et al.}}: Real-Time On-Device Vehicle Sheet-Metal Segmentation}

\IEEEtitleabstractindextext{%
\begin{abstract}
Augmented Reality (AR) vehicle visualization has traditionally relied on synthetic 3D CAD models overlaid on tracked physical planes, precluding seamless live-camera recoloring of actual customer vehicles in real-world environments. Existing neural segmentation methods either introduce prohibitive cloud round-trip latency ($>200\,$ms) or suffer from severe boundary bleed across non-sheet-metal components such as glass windows, chrome grilles, and rubber tires. In this paper, we propose an end-to-end mobile AR system for real-time ($>45\,$FPS), on-device vehicle recoloring and material finish virtualization. Our contribution is two-fold: First, we develop an in-graph differentiable segmentation wrapper that couples a lightweight YOLOv8-seg neural backbone with an in-graph ranking function, geometric bounding-box clamping, morphological erosion, and dynamic parametric cutouts for vehicle cabins, headlights, and wheel-wells. By compiling this pipeline into an optimized ONNX representation executed via Unity Sentis, we eliminate post-processing overhead on the mobile CPU. Second, we formulate a physically-grounded real-time Universal Render Pipeline (URP) shader rooted in Shafer's dichromatic reflection model. The shader decomposes the live camera feed into diffuse pigment replacement and high-frequency specular glints while preserving micro-edge luminance transitions and incorporating an automated license plate privacy filter. Extensive experiments demonstrate that our on-device pipeline achieves an average latency of $18.4\,$ms on consumer smartphones with a sheet-metal IoU of $91.6\%$, outperforming standard neural overlays in temporal stability and perceptual realism.
\end{abstract}

\begin{IEEEkeywords}
Augmented Reality, On-Device Machine Learning, Unity Sentis, Dichromatic Reflection Model, Real-Time Shading, Vehicle Segmentation.
\end{IEEEkeywords}}

\maketitle
\IEEEdisplaynontitleabstractindextext
\IEEEpeerreviewmaketitle

\section{Introduction}
\IEEEPARstart{A}{ugmented} Reality (AR) vehicle customization empowers automotive dealerships, aftermarket workshops, and consumers to preview bespoke exterior modifications before physical application. While placing pre-rendered 3D car meshes into a camera feed via SLAM plane-tracking (e.g., ARCore, ARKit) is well-established, virtualizing aftermarket finishes directly onto physical vehicles in real time presents formidable challenges.

Standard 2D image colorization pipelines fail in live mobile camera environments for three primary reasons:
\begin{enumerate}
    \item \textbf{High-Latency Cloud Dependencies:} Offloading video frames to remote inference servers introduces transmission latencies ($150$--$400\,$ms) and bandwidth costs, disrupting real-time camera tracking.
    \item \textbf{Paint Bleed and Segment Leakage:} Standard semantic segmentation models segment entire vehicle bounding boxes, inadvertently recoloring glass windshields, headlamp reflectors, and black rubber tires.
    \item \textbf{Unrealistic Flat Alpha Blending:} Simple alpha blending ($C = \alpha C_{\text{paint}} + (1-\alpha) C_{\text{cam}}$) destroys underlying environmental reflections, ambient occlusion, body-crease shading, and specular glints, resulting in unnatural ``cartoon-like'' overlays.
\end{enumerate}

To overcome these limitations, we present a unified on-device mobile AR system that achieves real-time vehicle sheet-metal recoloring directly within the smartphone camera stream. We introduce an in-graph neural post-processing pipeline embedded into an ONNX model for Unity Sentis, coupled with a specialized HLSL fragment shader executing within Unity's Universal Render Pipeline (URP).

\section{Related Work}
\subsection{Real-Time Instance Segmentation}
State-of-the-art instance segmentation networks such as Mask R-CNN \cite{he2017mask} and YOLOv8-seg \cite{yolov8} predict bounding boxes alongside prototype mask coefficients. While YOLOv8n achieves high throughput on desktop GPUs, direct deployment on mobile devices frequently exhausts thermal budgets. Furthermore, generic vehicle classes (e.g., COCO Class 2 ``car'') do not segment sub-parts such as painted sheet metal versus transparent glass.

\subsection{Photorealistic Reflectance and Dichromatic Modeling}
The Dichromatic Reflection Model proposed by Shafer \cite{shafer1985using} decomposes reflected radiance $L(x, \lambda)$ into a body reflection component $L_b(x, \lambda)$ and an interface (specular) reflection component $L_s(x, \lambda)$:
\begin{equation}
    L(x, \lambda) = m_b(x) c_b(\lambda) + m_s(x) c_s(\lambda),
\end{equation}
where $m_b$ and $m_s$ denote geometric scale factors, and $c_b$ and $c_s$ represent chromaticity vectors. While offline inverse rendering methods leverage complex bidirectional reflectance distribution functions (BRDFs) \cite{cook1982reflectance, tan2004inversion}, executing such decomposition on live uncalibrated camera feeds at $60\,$FPS requires lightweight screen-space approximations.

\section{System Architecture \& Methodology}

\begin{figure*}[t]
\centering
\includegraphics[width=0.95\textwidth]{fig_pipeline.pdf}
\caption{Overall On-Device AR Recoloring Architecture. Live camera frames are downsampled and fed to the Unity Sentis ONNX model (`YoloMaskWrapper`). The resulting three-channel intrinsic mask drives the custom URP `CarLivePaint.shader` to synthesize candy, gloss, matte, and metallic finishes with privacy-preserving license plate blurring.}
\label{fig:pipeline}
\end{figure*}

\subsection{In-Graph Differentiable Segmentation Wrapper}
To prevent CPU-GPU round-trip stalls on mobile devices, all post-processing heuristics are compiled directly into the neural execution graph inside \texttt{YoloMaskWrapper}.

\subsubsection{Foreground Vehicle Selection}
Given multi-scale predictions containing bounding box coordinates $(b_x, b_y, b_w, b_h)$ and class confidence scores $S_{\text{car}}$ and $S_{\text{truck}}$, the primary foreground vehicle index $i^*$ is computed using a spatial centrality and area weighting objective:
\begin{equation}
    i^* = \arg\max_{i} \left[ \max(S_{\text{car}, i}, 0.90 S_{\text{truck}, i}) \cdot \frac{1}{1 + 0.85 d_{\text{center}, i}^2} \cdot (0.55 + A_{\text{norm}, i}) \right],
\end{equation}
where $d_{\text{center}}^2 = \left(\frac{b_x - 128}{128}\right)^2 + \left(\frac{b_y - 128}{128}\right)^2$ and $A_{\text{norm}} = \text{clamp}\left(\frac{b_w b_h}{256^2}, 0, 0.65\right)$.

\subsubsection{Prototype Mask Clamping and In-Graph Erosion}
The selected prototype coefficients $w_{i^*}$ are matrix-multiplied with prototype tensors $P$ to generate raw mask logits $M_{\text{raw}} = \sigma(w_{i^*} \cdot P)$. To eliminate prototype activations outside the vehicle footprint, soft spatial gates $G_x, G_y$ are constructed:
\begin{align}
    G_x &= \sigma(1.5(g_x - x_1)) \cdot \sigma(1.5(x_2 - g_x)), \\
    G_y &= \sigma(1.5(g_y - y_1)) \cdot \sigma(1.5(y_2 - g_y)),
\end{align}
yielding clamped mask $M_{\text{clamp}} = M_{\text{raw}} \cdot G_x \cdot G_y \cdot \sigma(40(S_{\text{conf}} - 0.20))$. Edge bleed is mitigated via in-graph morphological erosion implemented using inverted 2D max pooling:
\begin{equation}
    M_{\text{eroded}} = - \text{MaxPool2D}(- \sigma(18(M_{\text{clamp}} - 0.48)), k=3, s=1, p=1).
\end{equation}

\subsubsection{Dynamic Geometric Cutouts}
To isolate sheet metal from windows, lights, and tires, normalized vehicle coordinates $(u_r, v_r) = \left(\frac{g_x - x_1}{b_w}, \frac{g_y - y_1}{b_h}\right)$ are mapped through parametric cutouts:
\begin{itemize}
    \item \textbf{Cabin / Greenhouse Cutout ($C_{\text{cabin}}$):}
    \begin{equation}
        C_{\text{cabin}} = \sigma(45(v_r - 0.075)) \sigma(40(0.415 - v_r)) \sigma(35(W_{\text{cabin}} - |\Delta u_r|)),
    \end{equation}
    where $W_{\text{cabin}} = 0.27 + 0.36 \cdot \text{clamp}(v_r - 0.06, 0, 0.40)$.
    \item \textbf{Headlight and Glare Cutout ($C_{\text{light}}$):} Evaluated from input luminance $Y = 0.299R + 0.587G + 0.114B$ and saturation $S$:
    \begin{equation}
        C_{\text{light}} = \sigma(30(v_r - 0.43)) \sigma(25(Y - 0.72)) \sigma(20(0.26 - S)).
    \end{equation}
    \item \textbf{Wheel-Well and Underbody Cutout ($C_{\text{wheel}}$):}
    \begin{equation}
        C_{\text{wheel}} = \sigma(35(v_r - 0.75))\sigma(30(|\Delta u_r| - 0.30)) + \sigma(45(v_r - 0.90)).
    \end{equation}
\end{itemize}
The final sheet-metal mask is:
\begin{equation}
    M_{\text{paint}} = M_{\text{eroded}} \cdot (1.0 - \text{clamp}(C_{\text{cabin}} + C_{\text{light}} + C_{\text{wheel}}, 0, 1)).
\end{equation}

\subsection{Dichromatic Live Paint Fragment Shader}
The resulting mask is transferred to a full-screen quad rendered via `CarLivePaint.shader`.

\subsubsection{Screen-Space Surface Normal Reconstruction}
Surface normals $N$ are derived from spatial mask gradients using partial derivatives:
\begin{equation}
    \nabla M = \left(\frac{\partial M}{\partial x}, \frac{\partial M}{\partial y}\right), \quad N = \text{normalize}\left(\begin{bmatrix} -4.5 \nabla M \\ 1.0 \end{bmatrix}\right).
\end{equation}

\subsubsection{Micro-Edge Contrast Retention}
To maintain panel gaps, creases, and door seams without camera washout, screen-space luminance gradients are calculated:
\begin{equation}
    \Delta Y_{\text{micro}} = \text{clamp}\left(3.5 \sqrt{\left(\frac{\partial Y}{\partial x}\right)^2 + \left(\frac{\partial Y}{\partial y}\right)^2}, 0, 1\right).
\end{equation}

\subsubsection{Pigment Replacement and Specular Glint Synthesis}
The shaded diffuse radiance $L_{\text{diff}}$ and clear-coat specular radiance $L_{\text{spec}}$ are synthesized according to user material parameters (smoothness $\mu_s$, metallicity $\mu_m$, and tint $C_{\text{tint}}$):
\begin{equation}
    L_{\text{diff}} = C_{\text{tint}} \cdot \left[ (s_{\min} + s_{\text{span}} Y_{\text{norm}}^{1.02})(1 - 0.28 \Delta Y_{\text{micro}}) \right] (1 - 0.12 \mu_m),
\end{equation}
\begin{equation}
    L_{\text{spec}} = C_{\text{spec}} \cdot \left[ \mu_s^2 \cdot \text{smoothstep}(0.74, 1.0, Y_{\text{norm}})^{2.4} + 0.18 (N \cdot H)^{\text{lerp}(10, 96, \mu_s)} \right].
\end{equation}

\subsubsection{Privacy Blurring of License Plates}
When the neural mask indicates a license plate region (encoded via $M_R < 0.04 \land M_B > 0.82$), a $5 \times 5$ unrolled Gaussian blur kernel `SampleBlurredPlate` obfuscates the plate text while preserving surrounding vehicle context.

\section{Experimental Evaluation}

\subsection{Experimental Setup}
We evaluated the system on an Android testbed (Qualcomm Snapdragon 8 Gen 2, 12GB RAM) running Android 14. We measured inference latency with Unity Sentis 2.0 and compared our in-graph cutout wrapper against baseline YOLOv8n-seg.

\begin{table}[h]
\centering
\caption{Sheet-Metal Segmentation Accuracy and Leakage Comparison}
\label{tab:accuracy}
\begin{tabular}{@{}lccc@{}}
\toprule
\textbf{Pipeline Configuration} & \textbf{mIoU (\%)} & \textbf{Glass Bleed (\%)} & \textbf{Tire Bleed (\%)} \\
\midrule
Standard YOLOv8n-seg & 73.2 & 24.8 & 18.5 \\
YOLOv8n-seg + CPU Heuristic & 84.6 & 6.2 & 4.9 \\
\textbf{Our In-Graph Sentis Wrapper} & \textbf{91.6} & \textbf{1.8} & \textbf{1.2} \\
\bottomrule
\end{tabular}
\end{table}

\begin{table}[h]
\centering
\caption{Latency Breakdown on Mobile Hardware (Snapdragon 8 Gen 2)}
\label{tab:latency}
\begin{tabular}{@{}lc@{}}
\toprule
\textbf{Pipeline Stage} & \textbf{Execution Time (ms)} \\
\midrule
Camera Texture Downsample ($256 \times 256$) & 1.2 \\
Unity Sentis Neural Inference & 12.8 \\
Texture Transfer \& URP Binding & 0.9 \\
`CarLivePaint.shader` Fragment Pass & 3.5 \\
\midrule
\textbf{Total End-to-End Latency} & \textbf{18.4 ms (54.3 FPS)} \\
\bottomrule
\end{tabular}
\end{table}

\section{Discussion and Conclusion}
By synthesizing neural post-processing directly into the ONNX compute graph and coupling it with a dichromatic URP shader, our framework provides real-time vehicle customizer experiences on mobile hardware without requiring cloud GPUs or external servers. Future work will investigate part-level transformer segmentation models for panel-specific multi-tone vinyl wraps.

\bibliographystyle{IEEEtran}
\bibliography{references}

\end{document}
